% ----- CHAUPAI 17 -----

\newpage

\subsection*{\deva{सप्तदश चौपाई} \(Seventeenth Chaupai\)}
\addcontentsline{toc}{subsection}{\deva{सप्तदश चौपाई} \(Seventeenth Chaupai\) — \deva{तुम्हरो मंत्र...}}

\begin{minipage}[t]{0.55\textwidth}
\vspace{0pt}

{\large \linespread{1.5}\selectfont
\deva{तुम्हरो मंत्र बिभीषन माना।}\\
\deva{लंकेस्वर भए सब जग जाना॥}
\par}

\vspace{0.5em}

{\large \linespread{1.5}\selectfont
\telu{తుమ్హరో మంత్ర బిభీషన మానా।}\\
\telu{లంకేశ్వర భఏ సబ జగ జానా॥}
\par}

\vspace{0.5em}

{\large \linespread{1.3}\selectfont
\textit{tumharo mantra bibhīṣana mānā |}\\
\textit{laṅkesvara bhae saba jaga jānā ||}
\par}
\end{minipage}%
\hfill
\begin{minipage}[t]{0.42\textwidth}
\vspace{0pt}
\begin{tcolorbox}[anchorbox]
\textbf{The Bridge of Advocacy:} In the *Valmiki Ramayana* (*Yuddha Kanda, 17:50-68*), when Vibhishana sought refuge, Sugriva and other leaders suspected him of being a spy. Hanuman was the sole advisor who delivered a brilliant, analytical speech to Lord Rama, advocating for Vibhishana's acceptance. Because Vibhishana was accepted and protected based on Hanuman's fair advocacy (this 'Mantra' or counsel to Rama), he was crowned King of Lanka (*Lankeshvara*). It shows how a true leader's strategic advocacy can lift others to greatness.
\vspace{0.5em}\newline Hanuman's brilliant, logical counsel on behalf of Vibhishana secured his sanctuary and his ultimate elevation to king.\newline\null\hfill\href{https://www.valmiki.iitk.ac.in/sloka?field_kanda_tid=6&language=dv&field_sarga_value=17}{\textit{Yuddha Kanda, 17:50-68}}
\end{tcolorbox}
\end{minipage}

\vspace{0.5em}
\subsubsection*{\textbf{Visual Rhythm Map:}}
\begin{table}[h!]
\centering
\renewcommand{\arraystretch}{1.2}
\setlength{\tabcolsep}{0pt}
\begin{tabularx}{\textwidth}{|*{16}{>{\centering\arraybackslash\hsize=1.0\hsize}X|}}
\hline
\multicolumn{4}{|>{\centering\arraybackslash\hsize=4.0\hsize}X|}{\textbf{\laghu{तु}\laghu{म्ह}\guru{रो}} \par (4)} &
\multicolumn{3}{>{\centering\arraybackslash\hsize=3.0\hsize}X|}{\textbf{\guru{मं}\laghu{त्र}} \par (3)} &
\multicolumn{5}{>{\centering\arraybackslash\hsize=5.0\hsize}X|}{\textbf{\laghu{बि}\guru{भी}\laghu{ष}\laghu{न}} \par (5)} &
\multicolumn{4}{>{\centering\arraybackslash\hsize=4.0\hsize}X|}{\textbf{\guru{मा}\guru{ना}} \par (4)} \\
\hline
\end{tabularx}
\vspace{-1.5em}
\end{table}

\begin{table}[h!]
\centering
\renewcommand{\arraystretch}{1.2}
\noindent\makebox[\textwidth][l]{%
  {%
  \setlength{\tabcolsep}{0pt}%
  \begin{tabularx}{1.0625\textwidth}{|*{17}{>{\centering\arraybackslash\hsize=1.0\hsize}X|}}
  \hline
  \multicolumn{6}{|>{\centering\arraybackslash\hsize=6.0\hsize}X|}{\textbf{\guru{लं}\guru{के}\laghu{स्व}\laghu{र}} \par (6)} &
  \multicolumn{3}{>{\centering\arraybackslash\hsize=3.0\hsize}X|}{\textbf{\laghu{भ}\guru{ए}} \par (3)} &
  \multicolumn{2}{>{\centering\arraybackslash\hsize=2.0\hsize}X|}{\textbf{\laghu{स}\laghu{ब}} \par (2)} &
  \multicolumn{2}{>{\centering\arraybackslash\hsize=2.0\hsize}X|}{\textbf{\laghu{ज}\laghu{ग}} \par (2)} &
  \multicolumn{4}{>{\centering\arraybackslash\hsize=4.0\hsize}X|}{\textbf{\guru{जा}\guru{ना}} \par (4)} \\
  \hline
  \end{tabularx}%
  }%
  \hspace{0.01\textwidth}%
  \begin{minipage}{0.1\textwidth}
  \vspace{-0.5em}
  \centering
  {\Huge \textbf{17}}
  \end{minipage}%
}
\vspace{-1.5em}
\end{table}

\vspace{0.5em}
\textbf{ \deva{सरल-संस्कृतम्}:}\\
 \deva{भवतः मन्त्रं (परामर्शं / उपदेशं) विभीषणः अङ्गीकृतवान्।}\\
 \deva{तेन सः लङ्केश्वरः (लङ्कायाः राजा) अभवत् इति सर्वं जगत् जानाति।} \\

Vibhishana heeded your counsel. 
By doing so, he became the King of Lanka, a fact the entire world knows.

\newpage

\vspace{1em}
\textbf{\deva{प्रतिपदार्थः} (Meaning of each word):}
\begin{table}[h!]
\centering
\renewcommand{\arraystretch}{1.2}
\begin{tabularx}{\textwidth}{@{} l l X X @{}}
\toprule
\textbf{अवधी} & \textbf{संस्कृतम्} & \textbf{English} & \textbf{Grammar \& Notes} \\
\midrule
\deva{तुम्हरो} / \telu{తుమ్హరో} / \textit{tumharo}  &  \deva{भवतः}  &  Your  &  Awadhi possessive; syncopated from \deva{तुम्हारो} \\
\deva{मंत्र} / \telu{మంత్ర} / \textit{mantra}  &  \deva{मन्त्रम् (उपदेशम्)}  &  Word / Counsel  &   \\
\deva{बिभीषन} / \telu{బిభీషన} / \textit{bibhīṣana}  &  \deva{विभीषणः}  &  Vibhishana  &  $v \rightarrow b$  \\
\deva{माना} / \telu{మానా} / \textit{mānā}  &  \deva{अङ्गीकृतवान्}  &  Accepted / Heeded  &   \\
\deva{लंकेस्वर} / \telu{లంకేశ్వర} / \textit{laṅkesvara}  &  \deva{लङ्केश्वरः}  &  King of Lanka  &  Compound: \deva{लंका} + \deva{ईश्वर} \\
\deva{भए} / \telu{భఏ} / \textit{bhae}  &  \deva{अभवत्}  &  Became  &   \\
\deva{सब जग} / \telu{సబ జగ} / \textit{saba jaga}  &  \deva{सर्वं जगत्}  &  The whole world  &   \\
\deva{जाना} / \telu{జానా} / \textit{jānā}  &  \deva{जानाति}  &  Knows  &   \\
\bottomrule
\end{tabularx}
\end{table}



\newpage
